\chapter{Modeling Mutation Purely: The State Monad}
\label{ch:monad}

\chapterorient{A solver does not compute one tensor; it \emph{evolves} a bundle of
state---the current iterate, the residual it is chasing, the count of how many
steps it has taken---across many passes of a loop. In a mutable language you would
keep that bundle in a struct and overwrite it. We cannot: every record in this book
is immutable (\cref{ch:records}), so each step must take the \emph{old} bundle and
return a \emph{new} one. Done by hand, threading that bundle through every call is
tedious clerical work, and clerical work is where bugs breed. This chapter
introduces the discipline that automates the threading and---this is the real
prize---makes the \emph{one place} the state actually changes stand out plainly on
the page. We call it the \emph{state monad}. We will not approach it as an abstract
algebraic structure; we will build it as what it is: careful bookkeeping for a
single evolving value.}

\section{The problem: threading one bundle by hand}
\label{sec:problem}

Recall the workhorse record from \cref{ch:records}: a \code{SimState}, the run-time
state a solve carries from step to step. It holds the current iterate, an iteration
counter, a convergence flag, perhaps a residual:

\begin{lstlisting}[style=l4style]
SimState = { x: Tensor[N], it: Int, converged: Bool, residual: Float }
\end{lstlisting}

\noindent A solve is a chain of steps, each producing the next \code{SimState} from
the current one---exactly the value-threading $\vs_0 \to \vs_1 \to \vs_2$ of
\cref{ch:bigidea}. Because records are immutable, ``produce the next state'' means
``read the old one and build a fresh one with the spread.'' A single step that
bumps the counter and records a new iterate looks like this:

\begin{lstlisting}[style=l4style]
step :: SimState -> SimState
step s =
  let x'  = improve s.x        in    -- compute a better iterate (a pure call)
  let r'  = residual_of x'     in    -- its residual          (a pure call)
  { ...s, x: x', it: s.it + 1, residual: r' }   -- the fresh state
\end{lstlisting}

\noindent So far, so clean: one step is a pure function \code{SimState -> SimState},
and we already know how to write it. The trouble starts when a step is not a single
expression but a \emph{sequence} of sub-operations, several of which each need to
read \emph{and} advance the state. Then the threading becomes manual, and the manual
threading is where the eye glazes over. Picture a step that must (i) read the current
iterate, (ii) advance it and bump the counter, (iii) read the just-advanced state to
compute a residual, and (iv) record that residual. Threaded by hand, every sub-step
must be handed the \emph{latest} state and must return the \emph{next} one, and the
programmer is personally responsible for never reusing a stale one:

\begin{lstlisting}[style=l4style]
step_by_hand :: SimState -> SimState
step_by_hand s0 =
  let s1 = { ...s0, x: improve s0.x }      in   -- s1 is s0 advanced
  let s2 = { ...s1, it: s1.it + 1 }        in   -- s2 is s1, counter bumped
  let s3 = { ...s2, residual: residual_of s2.x } in   -- s3 reads s2, records r
  s3
\end{lstlisting}

\noindent Read that carefully and you will see the hazard. There are four state
values in flight---\code{s0}, \code{s1}, \code{s2}, \code{s3}---and the program is
correct only because each line refers to the \emph{immediately preceding} one. Write
\lstinline[style=l4style]|s1.it| where you meant \lstinline[style=l4style]|s2.it|,
or thread \code{s0} into the residual line instead of \code{s2}, and the bug is
silent: the types still check (every \code{s} is a \code{SimState}), the program
still runs, and it computes the wrong thing. The numbers \code{s0}, \code{s1},
\code{s2}, \code{s3} are pure clerical overhead---they carry no meaning of their own;
they exist only to pass the baton from one line to the next. \Cref{fig:byhand} draws
the baton-passing we are about to automate away.

\begin{figure}[t]
\centering
\begin{tikzpicture}[font=\footnotesize, node distance=10mm]
  \node[data, minimum width=12mm] (s0) {\code{s0}};
  \node[opbox, right=9mm of s0] (o1) {improve};
  \node[data, minimum width=12mm, right=9mm of o1] (s1) {\code{s1}};
  \node[opbox, right=9mm of s1] (o2) {bump \code{it}};
  \node[data, minimum width=12mm, right=9mm of o2] (s2) {\code{s2}};
  \node[opbox, right=9mm of s2] (o3) {record \code{r}};
  \node[data, minimum width=12mm, right=9mm of o3] (s3) {\code{s3}};
  \draw[flow] (s0)--(o1); \draw[flow] (o1)--(s1);
  \draw[flow] (s1)--(o2); \draw[flow] (o2)--(s2);
  \draw[flow] (s2)--(o3); \draw[flow] (o3)--(s3);
  \node[font=\scriptsize\itshape, text=simRust, align=center] at (4.45,-1.1)
    {each baton must be handed to the \emph{next} line by hand---\\
     name the wrong one and the bug is silent};
\end{tikzpicture}
\caption[Threading state by hand]{The hand-threaded step of \cref{sec:problem}.
Four distinct state values \code{s0}\dots\code{s3} are passed line to line like a
relay baton; each operation reads the latest and emits the next. The names carry no
meaning---they are pure plumbing---and the program is correct only because every
line happens to reference the one just before it. The state monad of this chapter
makes the baton-passing implicit, so the plumbing names vanish.}
\label{fig:byhand}
\end{figure}

\begin{keyidea}
A solve evolves \emph{one} bundle of state through many steps. Immutability forbids
overwriting it, so each step must read the old bundle and return a fresh one---and
when a step has several sub-operations, threading that bundle by hand means inventing
a chain of meaningless names (\code{s0}, \code{s1}, \code{s2}, \dots) whose only job
is to carry the latest state to the next line. That bookkeeping is mechanical,
repetitive, and a quiet source of bugs. The \emph{state monad} automates it: you
write each sub-operation as if the state were simply ``there,'' and the monad threads
the right version between the lines for you.
\end{keyidea}

\section{\texttt{Solve} as a type: a recipe that threads a state}
\label{sec:solvetype}

Here is the central move, and we make it as concretely as we can. We give a name to
``a computation that threads a \code{SimState} and, along the way, yields some
result of type \code{a}.'' The name is \code{Solve}:

\begin{lstlisting}[style=l4style]
type Solve a = StateT SimState Identity a
\end{lstlisting}

\noindent Do not be alarmed by the right-hand side; we will unpack \code{StateT} and
\code{Identity} in a moment and then never need them again. Read the \emph{left} side
first, in plain words:

\begin{quote}\itshape
A value of type \code{Solve a} is a recipe that, when you give it a starting
\code{SimState}, produces a result of type \code{a} \emph{and} a new \code{SimState}.
\end{quote}

\noindent That is the entire idea. A \code{Solve a} is not a number, not a tensor,
not a record---it is a \emph{recipe} (a small function, if you like) waiting for a
state to run against. Run it against a state and two things come out: a value (the
\code{a}) and an updated state. Concretely, you can picture \code{Solve a} as
\emph{being} a function of this shape:

\begin{lstlisting}[style=l4style]
-- a Solve a is, under the hood, a function of this shape:
SimState -> (a, SimState)         -- give me a state; I return a value AND a new state
\end{lstlisting}

\noindent The fancy spelling \code{StateT SimState Identity a} is just the library
name for exactly that function shape---\code{StateT} reads ``state transformer over
the state type \code{SimState},'' and \code{Identity} is a do-nothing wrapper that
says ``no other effect is layered underneath; threading the state is all that
happens.'' We could have written the type as \code{State SimState a} and meant the
same thing; the \code{StateT \dots\ Identity} form is the spelling the calculus
inherits, and it commits us to nothing beyond the one-line reading above. From here
on, whenever you see \code{Solve a}, think only: \emph{a recipe that threads a
\code{SimState} and hands back an \code{a}}.

\begin{figure}[t]
\centering
\begin{tikzpicture}[font=\footnotesize]
  % the recipe box
  \node[stage, minimum width=34mm, minimum height=15mm, align=center] (box) at (0,0)
    {\code{Solve a}\\[1pt]\scriptsize\itshape a recipe};
  % input state on the left
  \node[data, minimum width=20mm, left=16mm of box] (sin) {\code{SimState} (in)};
  \draw[flow] (sin) -- (box);
  % two outputs on the right
  \node[product, minimum width=16mm] (val) at (4.4,0.9) {value : \code{a}};
  \node[data, minimum width=22mm] (sout) at (4.5,-0.9) {\code{SimState} (out)};
  \draw[flow] (box.east) -- (val.west);
  \draw[flow] (box.east) -- (sout.west);
  \node[font=\scriptsize\itshape, text=simGray, align=center] at (0,-1.55)
    {run it against a state\dots};
  \node[font=\scriptsize\itshape, text=simGray, align=center] at (4.5,-1.75)
    {\dots out comes a value \emph{and} a new state};
\end{tikzpicture}
\caption[\texttt{Solve a} as a state-threading recipe]{The plain reading of
\code{Solve a}. It is not a value sitting in memory; it is a \emph{recipe} that
consumes an incoming \code{SimState} and produces two things---a result of type
\code{a} and an outgoing \code{SimState}. The type \code{StateT SimState Identity a}
is the library spelling of exactly the function
\code{SimState -> (a, SimState)}; \code{Identity} merely says ``no other effect is
stacked underneath.'' Everything in this chapter is built from recipes of this one
shape and ways of gluing them together.}
\label{fig:solvetype}
\end{figure}

\begin{notationbox}
The name \code{Solve} is reused throughout the solver chapters, and the type
parameter \code{a} is the type of the \emph{result the recipe yields}, not the state
it threads. A \code{Solve Float} threads a \code{SimState} and yields a
\code{Float}; a \code{Solve Bool} threads the same \code{SimState} and yields a
\code{Bool}. The threaded state is \emph{always} a \code{SimState} (that is fixed in
the type \code{Solve a = StateT SimState Identity a}); only the yielded value varies.
A very common case is \code{Solve ()}---a recipe that threads the state but yields no
interesting value, written with the empty tuple \code{()} (``unit''). Such a recipe
exists \emph{only} for its effect on the state.
\end{notationbox}

\section{The primitives: \texttt{do}, \texttt{get}, \texttt{modify}, \texttt{return}}
\label{sec:primitives}

A recipe is useless until we can write one and run several in order. The state monad
gives us exactly four pieces of vocabulary, and they are all you need. We introduce
them one at a time, in the order you reach for them.

\paragraph{\code{do}: stack recipes top to bottom.} The keyword \code{do} sequences
recipes, top to bottom, exactly the way a list of imperative statements reads. Inside
a \code{do} block, each line is a recipe; the block runs them in order, threading the
state from each line into the next \emph{automatically}---this is the plumbing of
\cref{sec:problem}, now invisible. The block as a whole is itself a recipe (a bigger
\code{Solve}), so \code{do} blocks nest like any other expression.

\paragraph{\code{get}: read the current state.} The recipe \code{get} yields the
current \code{SimState} as its value and leaves the state unchanged. You bind its
result with the left-arrow \lstinline[style=l4style]|<-|:

\begin{lstlisting}[style=l4style]
do
  s <- get          -- s is now bound to the current SimState; state unchanged
  ...
\end{lstlisting}

\noindent After this line, the local name \code{s} \emph{is} the current state, and
you may dot into it (\lstinline[style=l4style]|s.it|, \lstinline[style=l4style]|s.x|)
like any record. Note the asymmetry with an ordinary \code{let}: a \code{let} binds
the value of a pure expression, while \lstinline[style=l4style]|s <- get| binds the
value \emph{a recipe yields}---here, the state itself. The left-arrow is how you reach
inside a recipe and name what it produced.

\paragraph{\code{modify f} / \code{put s}: write the state.} The recipe
\lstinline[style=l4style]|modify f| applies the pure function \code{f} to the current
state and makes the result the new state; it yields no interesting value (its type is
\code{Solve ()}). This is the \emph{single most important primitive in the chapter},
because \code{modify} is the one operation that \emph{changes} the threaded state.
The function \code{f} is an ordinary, pure \code{SimState -> SimState}---typically a
spread:

\begin{lstlisting}[style=l4style]
modify (\s -> { ...s, it: s.it + 1 })   -- advance the state: bump the counter
\end{lstlisting}

\noindent Read it: ``replace the current state by a fresh one, like it but with
\code{it} incremented.'' The sibling \lstinline[style=l4style]|put s'| is the blunt
version: it discards the current state and installs \code{s'} wholesale. We will
almost always prefer \code{modify}, because \code{modify} \emph{names the change}
(the spread says exactly which field moved) while \code{put} hides it inside a
fully-rebuilt record. Indeed
\lstinline[style=l4style]|modify f = do { s <- get; put (f s) }|---\code{modify} is
just ``get, apply \code{f}, put,'' bundled into one honest verb.

\paragraph{\code{return a}: yield a value, change nothing.} The recipe
\lstinline[style=l4style]|return a| yields the value \code{a} and leaves the state
untouched. It is how a \code{do} block produces a final result, and how a pure value
is lifted into a recipe so it can sit in a \code{do} block beside the real effects.
\lstinline[style=l4style]|return| is the ``do nothing to the state, just hand back
this value'' recipe.

\begin{notationbox}
The four primitives partition cleanly by what they touch.
\lstinline[style=l4style]|get| \emph{reads} the state (yields it, changes nothing);
\lstinline[style=l4style]|modify f| and \lstinline[style=l4style]|put s| \emph{write}
the state; \lstinline[style=l4style]|return a| touches the state not at all (it only
yields a value). Binding a recipe's yielded value uses the left-arrow
\lstinline[style=l4style]|x <- recipe|; binding a pure expression uses
\lstinline[style=l4style]|let x = expr|. The two binding forms living side by side in
one \code{do} block is the whole texture of monadic code, and \cref{sec:effectpoint}
makes the distinction between them load-bearing.
\end{notationbox}

\subsection{Walking a tiny \texttt{do} block, line by line}
\label{sec:walkdo}

Let us slow all the way down and watch the hidden state thread through a three-line
\code{do} block. The block reads the state, advances the iterate and counter, and
yields the new iteration number:

\begin{lstlisting}[style=l4style]
bump :: Solve Int
bump = do
  s <- get                                     -- (1) read the state into s
  modify (\s -> { ...s, x: improve s.x,        -- (2) write: advance x and it
                       it: s.it + 1 })
  s' <- get                                    -- (3) re-read the advanced state
  return s'.it                                 -- (4) yield the new iteration count
\end{lstlisting}

\noindent The block is a recipe of type \code{Solve Int}: run it against a starting
state and it yields an \code{Int} (the new counter) and leaves behind an advanced
state. \Cref{fig:walkdo} draws the hidden state threading between the lines. Trace it
against a concrete starting state
\lstinline[style=l4style]|s0 = { x: x0, it: 7, ... }|:

\begin{enumerate}[label=(\arabic*), leftmargin=*, nosep]
  \item \lstinline[style=l4style]|s <- get| yields the current state \code{s0} and
  binds it to \code{s}; the threaded state is still \code{s0}.
  \item \lstinline[style=l4style]|modify (...)| replaces the state by the spread
  applied to \code{s0}---a fresh \lstinline[style=l4style]|{ x: improve x0, it: 8, ... }|.
  The threaded state is now this new record; call it \code{s1}.
  \item \lstinline[style=l4style]|s' <- get| yields the \emph{current} state---which
  is now \code{s1}, the advanced one---and binds it to \code{s'}.
  \item \lstinline[style=l4style]|return s'.it| yields
  \lstinline[style=l4style]|s'.it|, which is $8$, and leaves the state as \code{s1}.
\end{enumerate}

\noindent The crucial thing to see is what happened \emph{between} the lines. Nowhere
did we write \code{s0}, \code{s1} as plumbing names; the \code{do} block threaded
them for us. Line (1)'s \code{get} saw the state before the \code{modify}; line (3)'s
\code{get} saw the state \emph{after} it. Same word, \code{get}, two different
values---because the \code{do} block silently advanced the threaded state across the
\code{modify} in between. That silent advance is precisely the baton-passing of
\cref{fig:byhand}, now performed by the monad instead of by you.

\begin{figure}[t]
\centering
\begin{tikzpicture}[font=\footnotesize]
  % the four code lines as a column on the left
  \node[anchor=west, font=\ttfamily\scriptsize] (l1) at (-6.2,1.65) {s <- get};
  \node[anchor=west, font=\ttfamily\scriptsize] (l2) at (-6.2,0.55) {modify (...)};
  \node[anchor=west, font=\ttfamily\scriptsize] (l3) at (-6.2,-0.55) {s' <- get};
  \node[anchor=west, font=\ttfamily\scriptsize] (l4) at (-6.2,-1.65) {return s'.it};
  % the threaded-state spine on the right
  \node[data, minimum width=16mm] (st0) at (0,1.65) {\code{s0}};
  \node[data, minimum width=16mm] (st1) at (0,0.55) {\code{s0}};
  \node[data, fill=simBgRust, draw=simRust, minimum width=16mm] (st2) at (0,-0.55) {\code{s1}};
  \node[data, fill=simBgRust, draw=simRust, minimum width=16mm] (st3) at (0,-1.65) {\code{s1}};
  % thread arrows down the spine
  \draw[flow] (st0)--(st1);
  \draw[flow] (st1)-- node[right=1mm,font=\scriptsize,text=simRust]{\textbf{the one write}} (st2);
  \draw[flow] (st2)--(st3);
  % which line reads what
  \draw[refedge] (l1.east) -- node[above,font=\scriptsize]{reads} (st0.west);
  \draw[refedge, simRust] (l2.east) -- node[above,font=\scriptsize,text=simRust]{writes} (st1.west);
  \draw[refedge] (l3.east) -- node[above,font=\scriptsize]{reads} (st2.west);
  \draw[refedge] (l4.east) -- node[above,font=\scriptsize]{ignores} (st3.west);
  \node[font=\scriptsize\itshape, text=simGray, align=center] at (0,2.45)
    {hidden threaded state (you never name it)};
\end{tikzpicture}
\caption[The hidden state threaded between \texttt{do} lines]{The \code{do} block of
\cref{sec:walkdo} with its \emph{hidden} threaded state drawn on the right. The four
source lines sit on the left; the state spine on the right shows what value the
threaded \code{SimState} holds as control passes each line. The first \code{get}
(line 1) sees \code{s0}; the lone \code{modify} (line 2, in rust) is the single point
where the state changes, advancing \code{s0} to \code{s1}; the second \code{get}
(line 3) therefore sees the \emph{advanced} \code{s1}. You wrote none of these spine
values---the \code{do} block threaded them. Two \code{get}s, one word, two different
answers, because one \code{modify} sat between them.}
\label{fig:walkdo}
\end{figure}

\subsection{Why this is equivalent to the hand-threaded version}
\label{sec:equiv}

It is worth proving to yourself, once, that the \code{do} block of \cref{sec:walkdo}
and a hand-threaded version compute exactly the same thing---that the monad is
genuinely just hiding plumbing and adds no magic. \Cref{wex:equiv} does this side by
side. The upshot, which we state now and earn there, is a small desugaring rule: a
\code{do} block is shorthand for a chain of explicit ``run this recipe against the
current state, name what it yields, thread the new state into the rest'' steps. The
monad is sugar; underneath, the baton still gets passed---you just no longer have to
name it.

\section{The single effect point: in the monad versus outside it}
\label{sec:effectpoint}

We now reach the heart of the chapter, and the reason the state monad earns its place
in a book about \emph{describing} algorithms clearly. The discipline is this:

\begin{quote}\itshape
Put in the monad \emph{only} the operations that read or write the threaded
\code{SimState}. Everything else---every operator application, every dot product,
every scrap of ephemeral workspace---is a pure \lstinline[style=l4style]|let|-binding
\emph{outside} the monad.
\end{quote}

\noindent The payoff is enormous and immediate: because only state-touching
operations live in the monad, the \emph{one place the state changes}---the lone
\code{modify}---stands out on the page, surrounded by pure \code{let}s that visibly do
not touch it. The monad turns ``where does this solver's state actually mutate?'' from
a question you answer by reading the whole loop into a question you answer by looking
for the word \code{modify}.

\subsection{A Krylov step: one \texttt{modify} amid a sea of pure \texttt{let}s}
\label{sec:krylovstep}

The cleanest illustration is a single step of a Krylov solver (the family
\cref{ch:operators} and the solver chapters build in full). A Krylov step does a
\emph{lot}: it applies the system operator to the current iterate, takes one or two
inner products, forms a scaled combination, updates a search direction, recomputes a
residual. Naively you might expect all of that to be ``in the monad''---it is, after
all, the step that advances the solve. The discipline says otherwise. Watch which
lines touch the \code{SimState}:

\begin{lstlisting}[style=l4style]
krylov_step :: OpParams -> SimState -> Solve ()
krylov_step op = \s_in -> do
  -- ===== OUTSIDE the monad: pure let-bindings, the actual numerical work =====
  let Ap   = apply op.A s_in.p          in   -- operator apply  (pure)
  let alpha = s_in.rdot / dot s_in.p Ap in   -- step length     (pure dot products)
  let x'    = axpy alpha s_in.p s_in.x  in   -- new iterate      (pure axpy)
  let r'    = axpy (-alpha) Ap s_in.r   in   -- new residual     (pure axpy)
  let rdot' = dot r' r'                 in   -- residual energy  (pure)
  let beta  = rdot' / s_in.rdot         in   -- direction coeff  (pure)
  let p'    = axpy beta s_in.p r'       in   -- new direction    (pure)
  -- ===== INSIDE the monad: the ONE state write =====
  modify (\s -> { ...s, x: x', r: r', p: p', rdot: rdot',
                       it: s.it + 1, converged: rdot' < op.tol })
\end{lstlisting}

\noindent Count the monadic operations: there is exactly \emph{one}, the closing
\code{modify}. Everything above it---the operator apply \code{Ap}, the inner products
\code{alpha} and \code{rdot'} and \code{beta}, the new iterate \code{x'}, residual
\code{r'}, and direction \code{p'}---is a pure \code{let}-binding. Those lines are the
algorithm's actual arithmetic, and \emph{none of them touches the state}. They read
fields out of \code{s\_in} (a plain value, captured once at the top) and compute fresh
tensors, exactly the pure-function vocabulary of \cref{ch:bigidea}. Only the final
\code{modify} folds the freshly computed quantities back into the threaded
\code{SimState}, advancing it by one step. \Cref{fig:partition} draws the partition.

\begin{figure}[t]
\centering
\begin{tikzpicture}[font=\footnotesize]
  % --- the "outside the monad" cloud of pure work ---
  \node[draw=simGray!55, fill=simBgGray, rounded corners=4pt, thick,
        minimum width=58mm, minimum height=30mm] (cloud) at (-2.6,0) {};
  \node[anchor=north, font=\scriptsize\bfseries, text=simGray]
    at ($(cloud.north)+(0,-1.5mm)$) {OUTSIDE the monad --- pure \code{let}s};
  \foreach \y/\lbl in {0.55/{Ap = apply A p}, -0.05/{alpha, beta = dot ...},
                       -0.65/{x', r', p' = axpy ...}}{
    \node[opbox, minimum width=42mm, font=\ttfamily\scriptsize] at (-2.6,\y) {\lbl};
  }
  \node[anchor=south, font=\scriptsize\itshape, text=simGray]
    at ($(cloud.south)+(0,1mm)$) {reads fields, computes fresh tensors; touches no state};
  % --- the single modify, inside ---
  \node[draw=simRust, fill=simBgRust, rounded corners=4pt, very thick,
        minimum width=34mm, minimum height=16mm, align=center] (mod) at (4.0,0)
    {\scriptsize\bfseries INSIDE the monad\\[2pt]\ttfamily\scriptsize modify (\textbackslash s -> \{...\})\\[2pt]\scriptsize\itshape the ONE state write};
  % the computed values flow into the modify
  \draw[flow] (cloud.east) -- node[above,font=\scriptsize]{\code{x', r', p', \dots}} (mod.west);
\end{tikzpicture}
\caption[In the monad versus outside it]{The partition of a Krylov step
(\cref{sec:krylovstep}). The grey cloud on the left is everything \emph{outside} the
monad: pure \code{let}-bindings that apply the operator, take inner products, and form
the new iterate, residual, and direction. They read fields off the incoming state
value and compute fresh tensors; not one of them touches the threaded \code{SimState}.
The rust box on the right is the \emph{single} monadic operation---one
\code{modify}---that folds those computed values back into the state, advancing it by
one step. The arithmetic is voluminous; the effect is a single point. That asymmetry,
made visible, is what the monad buys.}
\label{fig:partition}
\end{figure}

\begin{keyidea}
The state monad makes the single mutation point \emph{visible}. In a Krylov step the
only thing \emph{in} the monad is the state write---one \code{modify} that bumps the
counter and folds in the new iterate, residual, and direction. Everything else---the
operator applies, the inner products, the ephemeral workspace---is a pure
\code{let}-binding \emph{outside} the monad. The rule of thumb is exact: \textbf{if an
operation reads or writes \code{SimState}, it is in the monad; otherwise it is a pure
function call.} Because the discipline holds, you find every place a solver's state
changes by searching for one word, \code{modify}---and in a whole Krylov step there
is exactly one.
\end{keyidea}

\subsection{The rule of thumb, and why the counter is the canonical effect}
\label{sec:counter}

Apply the rule of thumb line by line to \cref{sec:krylovstep} and a striking fact
emerges. The genuinely \emph{state-bearing} field---the one whose entire reason for
living in \code{SimState} is that it must persist and be read across steps---is the
iteration counter \code{it}. The heavy tensors (\code{x}, \code{r}, \code{p}) are
recomputed each step from the previous step's values; they ride in the state for
convenience, but the quantity that genuinely \emph{accumulates}, that the loop's
predicate reads to decide whether to stop (\cref{ch:fold}), is the humble counter.
In the simplest possible step---one that does nothing but advance the loop---the
\emph{only} monadic effect is the counter bump:

\begin{lstlisting}[style=l4style]
modify (\s -> { ...s, it: s.it + 1 })   -- the canonical, irreducible state effect
\end{lstlisting}

\noindent This is the irreducible core of ``the solve mutated.'' We flagged it back
in \cref{ch:records}: a solve's externally visible state mutates at essentially a
\emph{single point}, the counter, while the heavy fields are recomputed elsewhere.
The monad is what cashes that observation in: it confines the mutation to one named
\code{modify} and pushes everything else out into pure \code{let}s, so the counter
bump---the one true effect---is the one thing left inside.

\begin{pitfall}
The cardinal sin of monadic code is \emph{over-monadifying}: dragging pure work
\emph{into} the monad that has no business there. It is tempting, having met
\code{do} and \code{get}, to write every line as a monadic step---to \code{get} the
state, compute \code{Ap = apply A s.p}, \code{put} a state with \code{Ap} stashed in
a scratch field, \code{get} again, and so on. Resist it utterly. An operator apply is
a pure function of its inputs; it does not read or write \code{SimState}, so it must
\emph{not} be in the monad. Putting it there does three bad things: it buries the one
real effect (the \code{modify}) in a crowd of fake ones, so the single mutation point
no longer stands out; it forces ephemeral workspace (\code{Ap}, \code{alpha}) into the
persistent state, where it does not belong (\cref{ch:strata}); and it defeats the
demand-driven pruning of \cref{ch:fold}, which can only eliminate \emph{pure}
unobserved work. The discipline is one-directional: pure work stays out; only
state reads and writes go in. When in doubt, ask the rule of thumb---\emph{does this
line read or write \code{SimState}?}---and if the answer is no, it is a \code{let}.
\end{pitfall}

\section{\texttt{execState}: the boundary where the monad evaporates}
\label{sec:execstate}

We have a way to build state-threading recipes and a discipline for keeping them
clean. One question remains: how do we \emph{run} one? A \code{Solve a} is a recipe,
not a value; to get an answer out we must supply a starting state and let the recipe
thread it through to the end. That is what \code{execState} does.

\begin{lstlisting}[style=l4style]
execState :: Solve a -> SimState -> SimState
\end{lstlisting}

\noindent Read the signature: \code{execState} takes a recipe (\code{Solve a}) and an
initial \code{SimState}, runs the recipe against it---threading the state through
every \code{modify} inside---and returns the \emph{final state}. (It is called
\code{exec} because it discards the recipe's yielded value \code{a} and keeps only the
final \emph{state}; a sibling \code{evalState} keeps the value and discards the state,
and \code{runState} keeps both. For a solve we want the final state, so
\code{execState} is our tool.) The entire monadic apparatus---the \code{do} blocks,
the \code{get}s and \code{modify}s, the hidden threading---is \emph{scaffolding that
exists only between} the initial state going in and the final state coming out.
Apply \code{execState} and the scaffolding evaporates: what is left is a plain
function from one \code{SimState} to another.

This is what lets the simulator stay pure on the outside. The top-level solve is
written:

\begin{lstlisting}[style=l4style]
solve :: OpParams -> Inputs -> SimState
solve op inp = execState (solve_loop op inp) (initial_state inp)
\end{lstlisting}

\noindent Look at the type of \code{solve}: \code{OpParams -> Inputs -> SimState}.
There is no \code{Solve}, no monad, no \code{StateT} anywhere in sight. Inside,
\code{solve\_loop op inp} is a \code{Solve ()} recipe---the whole iterative
computation, written monadically with its \code{do} blocks and \code{modify}s---and
\code{initial\_state inp} is the starting \code{SimState}. \code{execState} runs the
one against the other and hands back the final \code{SimState}. From outside,
\code{solve} is an ordinary pure function: feed it parameters and inputs, get a final
state. The monad lived entirely inside, did its bookkeeping, and was gone by the time
the answer crossed the boundary. \Cref{fig:execstate} draws the evaporation.

\begin{figure}[t]
\centering
\begin{tikzpicture}[font=\footnotesize]
  % the dashed monadic interior
  \node[draw=simViolet, fill=simBgViolet, rounded corners=4pt, densely dashed, thick,
        minimum width=46mm, minimum height=24mm] (mon) at (0,0) {};
  \node[anchor=north, font=\scriptsize\bfseries, text=simViolet]
    at ($(mon.north)+(0,-1.5mm)$) {inside: the \code{Solve} monad};
  \node[opbox, minimum width=34mm, font=\ttfamily\scriptsize] (a) at (0,0.45)
    {do \{ \dots\ modify \dots\ \}};
  \node[font=\scriptsize\itshape, text=simViolet, align=center] at (0,-0.55)
    {\code{get}s, \code{modify}s,\\hidden threading};
  % initial state in
  \node[data, minimum width=20mm] (s0) at (-5.0,0) {\code{initial\_state}};
  \draw[flow] (s0) -- (mon.west);
  % final state out
  \node[product, minimum width=20mm] (sf) at (5.0,0) {final \code{SimState}};
  \draw[flow] (mon.east) -- (sf.west);
  % the boundary label
  \node[font=\scriptsize\bfseries, text=simRust] at (-2.5,1.55) {\code{execState}};
  \node[font=\scriptsize\bfseries, text=simRust] at (2.5,1.55) {boundary};
  \node[font=\scriptsize\itshape, text=simGray, align=center] at (0,-1.95)
    {seen from outside: a pure \code{SimState -> SimState} --- the monad is gone};
\end{tikzpicture}
\caption[\texttt{execState} as the evaporation boundary]{\code{execState} runs a
\code{Solve} recipe against an initial state and returns the final state. The dashed
violet interior is the monadic computation---\code{do} blocks, \code{get}s,
\code{modify}s, the hidden threading of \cref{fig:walkdo}. \code{execState} is the
boundary: a \code{SimState} goes in on the left, the recipe threads it through every
internal effect, and a final \code{SimState} comes out on the right. From outside the
boundary there is no monad to be seen---\code{solve} has the plain type
\code{OpParams -> Inputs -> SimState}. The monad is scaffolding that does its work
inside and evaporates at the boundary, leaving the simulator pure on the outside.}
\label{fig:execstate}
\end{figure}

\begin{keyidea}
The monad is scaffolding that evaporates at the boundary. Inside \code{solve\_loop}
the state is threaded monadically, with \code{modify} marking each write; but
\code{execState (solve\_loop \dots) (initial\_state \dots)} runs that recipe against a
starting state and collapses the whole thing back into a \emph{pure} function
\code{SimState -> SimState}. The top-level \code{solve} has no monad in its type at
all. We get the local convenience of imperative-looking, automatically-threaded state
\emph{and} a globally pure simulator---the monad confined to the inside, gone by the
time any answer crosses out.
\end{keyidea}

\section{Why a monad and not just a fold}
\label{sec:whymonad}

A fair challenge: \cref{ch:fold} already gave us \code{iterate\_while}, which threads
a carry through a loop. If a solve is ``thread a \code{SimState} through steps until
converged,'' why not simply fold the \code{SimState} as the carry and be done? Why
introduce a monad at all?

The honest answer is that for the \emph{simplest} solvers you could. A non-restarted
solver with a single straight-line loop is perfectly well written as
\code{iterate\_while} with \code{SimState} as the carry. The monad earns its keep
when the control structure grows past a single flat loop, and three features push it
there.

\paragraph{Interior termination.} A real solver does not run a fixed number of steps;
it short-circuits the instant it detects convergence---and detection happens
\emph{mid-step}, on a residual the step just computed. Expressing ``compute the
residual, and if it is small enough, stop \emph{now}, before doing the rest'' is
exactly the kind of early exit a \code{do} block with a conditional makes plain at the
coordination layer. The fold's predicate, by contrast, fires only \emph{between}
steps (\cref{ch:fold}), so a pure fold must always finish a full step before it can
test; the monad lets the exit decision sit \emph{inside} the step where the residual
was born.

\paragraph{Inner and outer loop structure.} Restarted Krylov methods (GMRES is the
canonical case, built in \cref{ch:operators}) have an \emph{outer} cycle whose body is
itself an \emph{inner} loop. The inner loop reads \code{SimState.it} to honour the
global iteration cap and writes it back each step; the outer loop folds an accumulated
correction into \code{SimState.x} once per cycle. That is two-level read/write traffic
on one shared state, and the monad makes both levels explicit and the threading
automatic across the nesting. A single flat fold cannot express ``a loop whose step is
another loop'' without manually re-threading the carry through both.

\paragraph{Threading side-data.} The iterate \code{x} lives in \code{SimState}; the
Krylov basis lives in an ephemeral bundle reborn each restart; the correction step
\lstinline[style=l4style]|x += V*y| reads \emph{both}. The monad makes the moment the
\code{SimState} is touched a single named \code{modify}, while the ephemeral bundle
rides alongside as a plain \code{let}-bound value (it does not belong in the monad---it
does not persist across restarts). The monad cleanly separates ``the persistent state,
threaded'' from ``the per-cycle workspace, passed as a value.''

\noindent In every case the monad does the \emph{same} fundamental thing the fold
does---thread a value through steps---but it confines the threading of the
\emph{persistent} state to named \code{modify} points and lets the surrounding control
structure (early exit, nesting, side-data) be written plainly. And the connection to
\cref{ch:fold} is intact, not severed: the loop is \emph{still} \code{iterate\_while}.
The monad is simply what the carry threads \emph{through} when the carry is a
persistent simulator state and the control flow is richer than one flat loop. Fold and
monad are partners: the fold is the loop; the monad is the discipline for the state the
loop carries.

\begin{notationbox}
The three termination reasons a restarted solve can hit---converged on the residual
proxy, exhausted the global iteration cap, filled the per-cycle basis---collapse into
a single small sum type, \lstinline[style=l4style]@Outcome = Continue | Done Bool@.
The recipe \code{restart\_cycle} classifies the cycle's result once, at the boundary,
into an \code{Outcome}; \code{solve\_loop} pattern-matches on it and either recurses
or stops. This replaces a scatter of inner-loop break conditions with one decision
site---a clean payoff of writing the coordination monadically. The \code{Bool} inside
\code{Done} carries whether the stop was a convergence (\code{Done True}) or an
exhaustion (\code{Done False}), so the outer loop's final \code{modify} into
\code{SimState} is uniform.
\end{notationbox}

\section{Two faces of one name: the \texttt{Solve} effect versus the \texttt{Solve} record}
\label{sec:overload}

A reader who has been tracking the types carefully will have noticed something
slippery, and we must name it before it causes confusion. The word \code{Solve} is
\emph{overloaded}: it denotes two different things that happen to share a name.

\paragraph{\code{Solve} the effect.} This is the subject of the whole chapter so far:
\code{Solve a = StateT SimState Identity a}, the \emph{monad} that threads a
\code{SimState} and yields an \code{a}. It is a \emph{type constructor}---it takes a
type \code{a} (the yielded-value type) and produces the type of a recipe. ``\code{Solve}
the effect'' answers the question \emph{how is the state threaded?}

\paragraph{\code{Solve} the record.} When a single Krylov step is written out, the
\emph{value} it hands back is a record, and the calculus writes that record with the
same word: \lstinline[style=l4style]|Solve { sim, krylov, outputs }|. This is the
\code{SolveResult} bundle previewed in \cref{ch:records}: the next state (\code{sim}),
the next working bundle (\code{krylov}), and the demand-prunable per-step readouts
(\code{outputs}). ``\code{Solve} the record'' answers a different question: \emph{what
does one step yield?}

\noindent The two are genuinely distinct, and conflating them is the trap. The
\emph{effect} is the threading discipline; the \emph{record} is the shape of the value
yielded. They touch at exactly one field: the record's \code{sim} field is the one
that is discharged through the effect---it is the \code{SimState} the monad threads,
the product of the lone \code{modify}---while \code{krylov} and \code{outputs} are
plain returned values that ride alongside, untouched by the monad. So the rule of
thumb of \cref{sec:effectpoint} reappears at the level of the return record: of the
three fields the record hands back, \emph{one} (\code{sim}) is effectful and \emph{two}
(\code{krylov}, \code{outputs}) are pure. \Cref{fig:overload} disambiguates the two
faces.

\begin{figure}[t]
\centering
\begin{tikzpicture}[font=\footnotesize]
  % LEFT: Solve the effect
  \begin{scope}
    \node[font=\small\bfseries, text=simViolet] at (0,2.2) {\code{Solve} the \emph{effect}};
    \node[stage, minimum width=40mm, align=center] (eff) at (0,1.0)
      {\ttfamily\scriptsize Solve a = StateT SimState Identity a};
    \node[font=\scriptsize\itshape, text=simGray, align=center] at (0,0.0)
      {a type constructor:\\\emph{how} \code{SimState} is threaded};
    \node[font=\scriptsize, text=simInk, align=center] at (0,-1.0)
      {answers:\\``how is state threaded?''};
  \end{scope}
  \draw[simGray!50, thick] (3.4,-1.6) -- (3.4,2.4);
  % RIGHT: Solve the record
  \begin{scope}[xshift=7.0cm]
    \node[font=\small\bfseries, text=simBlue] at (0,2.2) {\code{Solve} the \emph{record}};
    \node[draw=simBlue, fill=simBgBlue, rounded corners=3pt, thick,
          minimum width=42mm, minimum height=18mm, inner sep=4pt] (rec) at (0,0.6) {};
    \node[anchor=north, font=\scriptsize\bfseries, text=simBlue] at ($(rec.north)+(0,-1mm)$)
      {\ttfamily Solve \{ ... \}};
    \node[anchor=west, font=\ttfamily\scriptsize, text=simRust]
      at ($(rec.north west)+(2.5mm,-6mm)$) {sim\ \ \itshape(effectful)};
    \node[anchor=west, font=\ttfamily\scriptsize]
      at ($(rec.north west)+(2.5mm,-10mm)$) {krylov, outputs\ \itshape(pure)};
    \node[font=\scriptsize, text=simInk, align=center] at (0,-1.0)
      {answers:\\``what does one step yield?''};
  \end{scope}
\end{tikzpicture}
\caption[The \texttt{Solve} overload disambiguated]{The two faces of \code{Solve}.
\emph{Left:} \code{Solve} the \emph{effect}, the monad
\code{StateT SimState Identity a}---a type constructor that answers \emph{how} the
\code{SimState} is threaded across the loop. \emph{Right:} \code{Solve} the
\emph{record}, the bundle \lstinline[style=l4style]|Solve { sim, krylov, outputs }| a
single step yields---a value that answers \emph{what} one step hands back. They meet
at one field: the record's \code{sim} (rust) is discharged \emph{through} the effect
(it is the threaded \code{SimState}), while \code{krylov} and \code{outputs} are plain
returned values the monad never touches. Same name, two questions; the \code{sim}
field is the bridge.}
\label{fig:overload}
\end{figure}

\begin{pitfall}
Do not read \lstinline[style=l4style]|Solve { sim, krylov, outputs }| as ``the monad
holds three pieces of state.'' It does not. The monad threads \emph{one} thing, a
\code{SimState}, and that one thing is the record's \code{sim} field alone. The
\code{krylov} and \code{outputs} fields are ordinary values the step returns
\emph{beside} the monadic effect---the working bundle (which is reborn each restart and
so deliberately is \emph{not} monadic state, lest its lifetime be misrepresented) and
the prunable readouts. When you see the word \code{Solve} in a signature, check which
face you are looking at: a bare \code{Solve a} (or \code{Solve ()}) is the effect; a
\lstinline[style=l4style]|Solve { ... }| with braces is the return record, and only its
\code{sim} field rides the effect.
\end{pitfall}

\section{A physical reprise: stream, then collide, threaded purely}
\label{sec:lbm}

To ground the whole chapter in something you can picture, we return to the
lattice-Boltzmann step we met in \cref{ch:bigidea}---stream, then collide---and show
it as the cleanest possible example of ``mutation modeled purely.'' It is the perfect
closing example precisely because, in its minimal form, it needs \emph{no} monad at
all: it is a pure \code{FluidState -> FluidState} function threaded by the plain fold
\code{iterate\_while\_pure} of \cref{ch:fold}. Seeing where the monad is \emph{not}
needed sharpens the sense of where it is.

Recall the rhythm. A lattice-Boltzmann fluid stores, at each cell of a grid, a little
tensor of \emph{populations}---how much fluid moves in each discrete direction. One
time step \emph{streams} (every population slides to the neighbor in its direction)
and then \emph{collides} (the populations in each cell relax toward local
equilibrium). In a classical code this is a nest of loops overwriting a population
array in place, often with an explicit scratch copy to avoid clobbering data a later
cell still needs:

\begin{lstlisting}[style=l4style, language=]
// classical C-style: overwrite in place, with a defensive copy
clone     = copy(current);              // guard against overwriting live data
streamed  = stream(clone);              // slide populations to neighbors
current   = collide(streamed);          // relax toward equilibrium, overwrite
\end{lstlisting}

\noindent The \lstinline[style=l4style, language=]|clone| is the tell. In the mutable
world you must defensively copy \code{current} before streaming, because streaming
reads a cell's neighbors and writing the result back in place would corrupt
neighbors not yet processed. The copy is pure ceremony forced by mutation. Now watch
the same step in our world, as a pure function building a fresh state:

\begin{lstlisting}[style=l4style]
step :: FluidState -> FluidState
step s =
  let streamed = stream  s.populations in    -- read s, produce a fresh tensor
  let collided = collide streamed       in    -- read it, produce another
  { ...s, populations: collided, t: s.t + 1 } -- a FRESH state; s untouched
\end{lstlisting}

\noindent There is no \code{clone}. There \emph{cannot} be a clobbering hazard,
because \code{stream} reads the immutable \code{s.populations} and \emph{returns} a
fresh tensor; the input it read can never change underneath it (\cref{ch:records}).
The defensive copy of the mutable version has \emph{evaporated}---in the immutable
world it was always a no-op in meaning, a workaround for a hazard that does not exist
here. That evaporation is the whole moral of \cref{ch:bigidea} made concrete: model
the mutation purely and the bookkeeping that mutation forced on you simply
disappears. \Cref{fig:lbm} traces a few steps.

\begin{figure}[t]
\centering
\begin{tikzpicture}[font=\footnotesize, node distance=8mm]
  % carry of FluidStates threaded by iterate_while_pure
  \node[data, minimum width=15mm, align=center] (f0) {\code{s}\(_0\)\\\scriptsize\(t=0\)};
  \node[opbox, right=10mm of f0, align=center] (st1) {stream\\then collide};
  \node[data, minimum width=15mm, align=center, right=10mm of st1] (f1) {\code{s}\(_1\)\\\scriptsize\(t=1\)};
  \node[opbox, right=10mm of f1, align=center] (st2) {stream\\then collide};
  \node[data, minimum width=15mm, align=center, right=10mm of st2] (f2) {\code{s}\(_2\)\\\scriptsize\(t=2\)};
  \node[font=\scriptsize, simGray, right=6mm of f2] (more) {$\cdots$};
  \draw[flow] (f0)--(st1); \draw[flow] (st1)-- node[above,font=\scriptsize]{fresh} (f1);
  \draw[flow] (f1)--(st2); \draw[flow] (st2)-- node[above,font=\scriptsize]{fresh} (f2);
  \draw[flow] (f2)--(more);
  % the evaporated clone
  \node[draw=simGray!45, densely dashed, fill=white, rounded corners=2pt,
        font=\scriptsize\itshape, text=simGray!70, align=center] (clone) at (1.7,-1.5)
    {\code{clone()}\\(no-op --- evaporates)};
  \draw[refedge, simGray!50] (clone.north) -- (st1.south);
  \node[font=\scriptsize\itshape, text=simGray, align=center] at (8.0,-1.5)
    {threaded by\\\code{iterate\_while\_pure}\\(no monad needed)};
\end{tikzpicture}
\caption[The LBM step threaded purely]{The lattice-Boltzmann step as a pure
\code{FluidState -> FluidState} function threaded by \code{iterate\_while\_pure}
(\cref{ch:fold}). Each step streams then collides the current populations and returns
a \emph{fresh} state with the time counter advanced; the chain
\code{s}\(_0 \to\) \code{s}\(_1 \to\) \code{s}\(_2\) is the value-threading of
\cref{ch:bigidea}. The defensive \code{clone()} of the mutable version (dashed, below)
is a \emph{no-op} here---it guarded against an in-place clobber that immutability makes
impossible, so it simply evaporates. Because the minimal step has no interior
termination, no nesting, and no persistent side-data beyond the carry itself, it needs
\emph{no monad}: the plain fold suffices.}
\label{fig:lbm}
\end{figure}

The lattice-Boltzmann step is the lower bound of the chapter's machinery: pure step,
plain fold, no monad. A Krylov solve sits at the upper bound: the same fold drives it
(\cref{ch:fold}), but now the carry is a persistent \code{SimState} with interior
termination and nesting, so the monad of this chapter manages the state the fold
carries, confining its one true mutation to a single \code{modify}. Between those two
poles live all the solvers of the book---and every one of them is the value-threading
of \cref{ch:bigidea}, dressed in exactly as much monadic bookkeeping as its control
structure demands, and no more.

\section{Worked examples}
\label{sec:worked}

\begin{workedexample}{One step, hand-threaded versus monadic---and they agree}{equiv}
We take a tiny two-effect computation and write it twice: once threading the state by
hand (passing \code{s} explicitly), once with the monad. Then we check that they
compute the same final state, so the monad is revealed as pure plumbing-hiding.

\textbf{The task.} From a starting state, (i) advance the iterate and bump the
counter, then (ii) record the residual of the advanced iterate. Two operations, each
touching the state.

\textbf{By hand.} Thread \code{s} explicitly, inventing names for each version:
\begin{lstlisting}[style=l4style]
by_hand :: SimState -> SimState
by_hand s0 =
  let s1 = { ...s0, x: improve s0.x, it: s0.it + 1 } in   -- (i) advance + bump
  let s2 = { ...s1, residual: residual_of s1.x }    in   -- (ii) read s1, record r
  s2
\end{lstlisting}
The correctness hinges on line (ii) reading \code{s1} (the advanced state), not
\code{s0}. The name \code{s1} is pure plumbing.

\textbf{With the monad.} Write each operation as a monadic step; the \code{do} block
threads the state:
\begin{lstlisting}[style=l4style]
monadic :: Solve ()
monadic = do
  modify (\s -> { ...s, x: improve s.x, it: s.it + 1 })   -- (i) advance + bump
  s' <- get                                               -- re-read the advanced state
  modify (\s -> { ...s, residual: residual_of s'.x })     -- (ii) record r
\end{lstlisting}

\tcblower
\textbf{Running the monadic version.} Apply \code{execState monadic s0} and thread by
the rules of \cref{sec:primitives}. The first \code{modify} replaces \code{s0} by
\lstinline[style=l4style]|{ ...s0, x: improve s0.x, it: s0.it + 1 }|---call it
\code{s1}, exactly the \code{s1} of the hand-threaded version. The
\lstinline[style=l4style]|s' <- get| binds \code{s'} to the \emph{current} state,
which is now \code{s1}. The second \code{modify} replaces \code{s1} by
\lstinline[style=l4style]|{ ...s1, residual: residual_of s'.x }|, and since
\code{s'} \emph{is} \code{s1}, that is
\lstinline[style=l4style]|{ ...s1, residual: residual_of s1.x }|---exactly the
hand-threaded \code{s2}. \code{execState} returns this final state.

\textbf{They agree.} The two final states are identical, character for character. The
difference is only that the hand-threaded version \emph{named} the intermediate
\code{s1} and had to reference it correctly by hand, while the monadic version let the
\code{do} block thread \code{s1} from the first \code{modify} into the \code{get} that
follows it. This is the desugaring promised in \cref{sec:equiv}: a \code{do} block is
shorthand for the explicit baton-passing, with the baton names supplied by the monad
instead of by you. The monad is sugar; the computation underneath is the value
threading of \cref{ch:bigidea}.
\end{workedexample}

\begin{workedexample}{The LBM step as a pure \texttt{FluidState -> FluidState}}{lbm}
We trace two lattice-Boltzmann steps numerically (schematically), watching the
\code{clone()} of the mutable version evaporate and the state thread as fresh
immutable values through \code{iterate\_while\_pure}.

\textbf{Setup.} A \code{FluidState} carries the population field and a time counter:
\begin{lstlisting}[style=l4style]
FluidState = { populations: Tensor[H, W, Q], t: Int }
\end{lstlisting}
where \code{Q} is the number of discrete velocity directions. The step streams then
collides; the driver folds it with the plain combinator:
\begin{lstlisting}[style=l4style]
run :: FluidState -> Int -> FluidState
run s0 nsteps =
  iterate_while_pure
    s0                              -- initial carry: the fluid state
    (\s -> s.t < nsteps)            -- predicate: pure on the carry (the time count)
    (\s -> let streamed = stream  s.populations in
           let collided = collide streamed       in
           { ...s, populations: collided, t: s.t + 1 })   -- fresh state, no clone
\end{lstlisting}

\tcblower
\textbf{The trace.} Take a schematic field where \code{stream} shifts populations and
\code{collide} nudges them toward equilibrium; we track only the time counter and the
\emph{identity} of the population tensor, to make the threading visible:
\begin{center}
\setlength{\tabcolsep}{8pt}
\renewcommand{\arraystretch}{1.15}
\begin{tabular}{@{}clll@{}}
\toprule
step & carry \code{t} & \code{streamed} (fresh) & result \code{populations} (fresh) \\
\midrule
0 & $0$ & $\mathrm{stream}(P_0)$ & $P_1 = \mathrm{collide}(\mathrm{stream}(P_0))$ \\
1 & $1$ & $\mathrm{stream}(P_1)$ & $P_2 = \mathrm{collide}(\mathrm{stream}(P_1))$ \\
\bottomrule
\end{tabular}
\end{center}
Each row reads the \emph{previous} population tensor ($P_0$, then $P_1$) and produces a
\emph{fresh} one ($P_1$, then $P_2$); the originals $P_0$, $P_1$ are never altered, so
no defensive copy is needed before streaming. In the mutable code,
\lstinline[style=l4style, language=]|clone = copy(current)| would sit at the top of
each step to protect $P_0$ from being clobbered while its neighbors are streamed; here
that copy is a \emph{no-op}---streaming an immutable $P_0$ cannot clobber it---so it
evaporates entirely.

\textbf{Why no monad.} Examine the step against the three pressures of
\cref{sec:whymonad}. There is no \emph{interior termination}: the loop runs a fixed
\code{nsteps}, tested between steps by the predicate. There is no \emph{inner/outer
nesting}: one flat loop. There is no persistent \emph{side-data} beyond the carry: the
\code{streamed} intermediate lives and dies inside one step as an ephemeral
\code{let}. All three pressures are absent, so the carry \emph{is} the entire state
and a plain fold suffices---no \code{Solve}, no \code{modify}, no \code{execState}.
This is the cleanest face of ``mutation modeled purely'': the same value-threading as
a Krylov solve, with none of the monadic bookkeeping, because the control structure
does not demand it. The monad is a tool you reach for exactly when the loop outgrows
this shape, and not a moment sooner.
\end{workedexample}

\summarysection
A solver evolves one bundle of state---the \code{SimState} holding the iterate,
counter, residual, and convergence flag---across many steps. Immutability forbids
overwriting it, so each step reads the old bundle and returns a fresh one; threading
that bundle by hand means inventing a chain of meaningless plumbing names (\code{s0},
\code{s1}, \dots) and is a quiet source of bugs (\cref{sec:problem}). The
\emph{state monad} automates the threading. A \code{Solve a = StateT SimState
Identity a} is a \emph{recipe} that, given a starting \code{SimState}, yields an
\code{a} and a new \code{SimState}---read plainly, the function
\code{SimState -> (a, SimState)} (\cref{sec:solvetype}). Four primitives suffice:
\code{do} stacks recipes top to bottom and threads the state between them; \code{get}
reads the current state; \code{modify f} / \code{put s} write it; \code{return a}
yields a value, changing nothing (\cref{sec:primitives}). The load-bearing discipline
is the \emph{single effect point}: only operations that read or write \code{SimState}
go in the monad; everything else---operator applies, inner products, ephemeral
workspace---is a pure \code{let} outside it. In a Krylov step the \emph{only} monadic
operation is one \code{modify}, surrounded by pure \code{let}s; the irreducible core of
that effect is the counter bump
\lstinline[style=l4style]|modify (\s -> { ...s, it: s.it + 1 })|
(\cref{sec:effectpoint},~\cref{sec:counter}). Over-monadifying---dragging pure work
into the monad---buries the one real effect and defeats pruning; the rule of thumb
guards against it. \code{execState} runs a recipe against an initial state and collapses
the whole monad back into a pure \code{SimState -> SimState}: the scaffolding
evaporates at the boundary, leaving \code{solve} pure on the outside
(\cref{sec:execstate}). The monad earns its keep over a plain fold when control flow
grows past one flat loop---interior termination, inner/outer nesting, side-data
threading---but the loop is \emph{still} \code{iterate\_while}; the monad is what the
carry threads through (\cref{sec:whymonad}). The name \code{Solve} is overloaded: the
\emph{effect} \code{StateT SimState Identity a} (how the state is threaded) versus the
\emph{record} \lstinline[style=l4style]|Solve { sim, krylov, outputs }| (what a step
yields), meeting at the one effectful field \code{sim} (\cref{sec:overload}). The
lattice-Boltzmann step is the chapter's lower bound: a pure
\code{FluidState -> FluidState} threaded by \code{iterate\_while\_pure}, with the
mutable world's defensive \code{clone()} evaporating into a no-op---mutation modeled
purely, needing no monad at all (\cref{sec:lbm}).

\furtherreading
The state monad is the classic functional-programming answer to ``how do I thread
mutable-looking state through pure code,'' and the canonical, readable introduction is
Wadler's tutorial on how monads tame state and effects~\citep{wadler1995monads}, which
develops exactly the \code{get}/\code{put}/\code{return} vocabulary used here, from the
ground up and with no category theory. The state monad and the \code{StateT}
transformer (our \code{StateT SimState Identity}) are standard equipment in the
language Peyton Jones documents~\citep{peytonjones2003}; readers who want the precise
desugaring of \code{do}-notation into \texttt{>>=} will find it there. We have used the
monad here only as disciplined bookkeeping for one evolving value; the solver chapters
(\cref{ch:operators} and the Krylov solvers that follow it) put it to work, and
\cref{ch:strata} explains \emph{why} the threaded \code{SimState} splits cleanly from
the build-time \code{OpParams} and the per-step \code{SolveResult}---the stratification
that makes the single effect point possible.

\exercisessection
\begin{enumerate}[nosep]
  \item In the hand-threaded \code{step\_by\_hand} of \cref{sec:problem}, suppose a
  typo made the third line read \lstinline[style=l4style]|residual_of s1.x| instead of
  \lstinline[style=l4style]|s2.x|. Does the program still type-check? Does it compute
  the right residual? Use this to explain, in one sentence, what class of bug the state
  monad eliminates.
  \item Read the type \code{Solve a = StateT SimState Identity a} aloud as a plain
  English sentence, without using the words ``monad'' or ``StateT.'' Then say what
  \code{Solve ()} and \code{Solve Bool} each yield, and what state each threads.
  \item For each of the four primitives \code{get}, \code{modify f}, \code{put s},
  \code{return a}, state whether it \emph{reads} the state, \emph{writes} the state, or
  \emph{neither}, and what value (if any) it yields. Which one is the only one that can
  change the threaded state?
  \item Walk the \code{do} block \code{bump} of \cref{sec:walkdo} against a starting
  state \lstinline[style=l4style]|{ x: x0, it: 3 }|. What value does each \code{get}
  yield? What does the block return? Explain why the two \code{get}s yield states with
  \emph{different} \code{it} fields even though they are the same word.
  \item Classify every line of the Krylov step in \cref{sec:krylovstep} as ``in the
  monad'' or ``outside the monad'' by applying the rule of thumb. How many lines are in
  the monad? Now argue why moving the line \lstinline[style=l4style]|let Ap = apply op.A s_in.p|
  \emph{into} the monad (via a \code{modify} that stashes \code{Ap} in a scratch field)
  would be a mistake---name two distinct things it breaks.
  \item The top-level \code{solve} has type \code{OpParams -> Inputs -> SimState}, with
  no \code{Solve} in sight, yet its body \code{solve\_loop} is a \code{Solve ()} recipe.
  Explain how \code{execState} reconciles these---what goes in, what comes out, and why
  the monad is invisible from outside.
  \item Give one concrete reason from \cref{sec:whymonad} that a restarted GMRES solve
  is awkward to write as a single flat \code{iterate\_while} with \code{SimState} as the
  carry. Then explain in what sense the loop is ``still \code{iterate\_while}'' even when
  the state is threaded monadically.
  \item Distinguish the two faces of \code{Solve}. Given the signature
  \lstinline[style=l4style]|krylov_step :: OpParams -> SimState -> Solve { sim, krylov, outputs }|,
  identify which \code{Solve} is the effect and which is the record, and say which
  \emph{single} field of the record is discharged through the effect and which two are
  plain returned values.
  \item In the LBM step of \cref{sec:lbm} the mutable version needs a defensive
  \code{clone()} but the pure version does not. Explain precisely what hazard the
  \code{clone()} guards against in the mutable world, and why immutability
  (\cref{ch:records}) makes that hazard impossible---so the copy becomes a no-op. Then
  state the three pressures of \cref{sec:whymonad} and confirm the minimal LBM step
  exhibits none, which is why it needs no monad.
\end{enumerate}
